\documentclass[a4,11pt]{aleph-notas}
% Para obtener solo el pdf, compilar con pdfLaTeX.
% Para obtener el xml compilar con XeLaTeX.

% -- Paquetes adicionales
\usepackage[handout]{aleph-moodle}
% \usepackage{aleph-moodle}
\moodleregisternewcommands
% Todos los comandos nuevos deben ir luego del comando anterior
\usepackage{aleph-comandos}


% -- Datos
\institucion{Facultad de Ciencias Exactas, Naturales y Ambientales}
\carrera{Catálogo STEM}
\asignatura{Matemática III}
\tema{Cuestionario No. 3: Integrales en varias variables}
\autor{Andrés Merino}
\fecha{Periodo 2026-1}

\logouno[0.16\textwidth]{Logos/logoPUCE_04_ac}
\definecolor{colortext}{HTML}{1957d1}
\definecolor{colordef}{HTML}{1957d1}
\fuente{montserrat}


\begin{document}

\encabezado

\vspace*{-8mm}
\tableofcontents

%%%%%%%%%%%%%%%%%%%%%%%%%%%%%%%%%%%%%%%%
\section{Indicaciones}
%%%%%%%%%%%%%%%%%%%%%%%%%%%%%%%%%%%%%%%%

Se plantean bancos de preguntas orientados a evaluar el \textbf{criterio}: «Interpreta el concepto de integrales en varias variables y reconocer sus elementos y propiedades», correspondiente al \textbf{resultado de aprendizaje}: Identificar los conceptos, teoremas y propiedades de funciones vectoriales, derivadas parciales e integral vectorial.

%%%%%%%%%%%%%%%%%%%%%%%%%%%%%%%%%%%%%%%%
\section{Banco de preguntas}
%%%%%%%%%%%%%%%%%%%%%%%%%%%%%%%%%%%%%%%%


%%%%%%%%%%%%%%%%%%%%%%%%%%%%%%%%%%%%%%%%
\begin{quiz}{Integral doble: concepto, elementos e interpretación}
%%%%%%%%%%%%%%%%%%%%%%%%%%%%%%%%%%%%%%%%

\begin{multi}[]%
    {ID-Concepto-01}
    Sea $f:\R^2\to\R$ un campo escalar continuo sobre una región $R$. El resultado de $\displaystyle\iint_R f\,dA$ es:
    \item* Un número real
    \item Un campo escalar
    \item Un campo vectorial
    \item Una función de dos variables
\end{multi}

\begin{multi}[]%
    {ID-Concepto-01-v1}
    Al evaluar la integral doble $\displaystyle\iint_R f\,dA$ de un campo escalar continuo $f$ sobre $R$, se obtiene:
    \item* Un número real
    \item Una función de dos variables
    \item Un campo vectorial
    \item Una región del plano
\end{multi}

\begin{multi}[]%
    {ID-Concepto-01-v2}
    ¿Qué tipo de objeto matemático es $\displaystyle\iint_R f\,dA$ cuando $f:\R^2\to\R$ es continua sobre $R$?
    \item* Un escalar (número real)
    \item Un vector
    \item Una superficie
    \item Una curva en el plano
\end{multi}

\begin{multi}[]%
    {ID-Concepto-01-v3}
    El cálculo de la integral doble de una función $f:\R^2\to\R$ sobre una región $R$ produce como resultado:
    \item* Un número real
    \item Otra función de $x$ e $y$
    \item Un campo vectorial
    \item Un plano tangente
\end{multi}

\begin{multi}[]%
    {ID-Concepto-01b}
    La integral doble $\displaystyle\iint_R f\,dA$ se define sobre:
    \item* Una región del plano $\R^2$
    \item Una curva del plano
    \item Un sólido de $\R^3$
    \item Un intervalo de $\R$
\end{multi}

\begin{multi}[]%
    {ID-Concepto-01b-v1}
    El dominio de integración de una integral doble $\displaystyle\iint_R f\,dA$ es:
    \item* Una región bidimensional del plano
    \item Un segmento de recta
    \item Un volumen de $\R^3$
    \item Un punto del plano
\end{multi}

\begin{multi}[]%
    {ID-Concepto-01b-v2}
    ¿Sobre qué tipo de conjunto se integra en una integral doble?
    \item* Una región del plano $\R^2$
    \item Una curva de $\R^2$
    \item Una superficie de $\R^3$
    \item Un intervalo cerrado $[a,b]$
\end{multi}

\begin{multi}[]%
    {ID-Concepto-01b-v3}
    En la expresión $\displaystyle\iint_R f\,dA$, el conjunto $R$ representa:
    \item* Una región plana (subconjunto de $\R^2$)
    \item Una trayectoria en el plano
    \item Un sólido tridimensional
    \item Un único número real
\end{multi}

\begin{multi}[]%
    {ID-Interp-01}
    Si $f(x,y)\geq 0$ sobre $R$, la integral $\displaystyle\iint_R f\,dA$ representa geométricamente:
    \item* El volumen del sólido bajo la gráfica de $f$ y sobre la región $R$.
    \item El área de la región $R$.
    \item La pendiente del plano tangente a la gráfica de $f$.
    \item La longitud de la frontera de $R$.
\end{multi}

\begin{multi}[]%
    {ID-Interp-01-v1}
    Para $f(x,y)\geq 0$ en $R$, ¿qué cantidad geométrica calcula $\displaystyle\iint_R f\,dA$?
    \item* El volumen comprendido entre la región $R$ y la superficie $z=f(x,y)$.
    \item El área de la región $R$.
    \item El perímetro de $R$.
    \item La altura máxima de la gráfica de $f$.
\end{multi}

\begin{multi}[]%
    {ID-Interp-01-v2}
    Cuando $f$ es no negativa sobre $R$, el valor de $\displaystyle\iint_R f\,dA$ coincide con:
    \item* El volumen del sólido limitado superiormente por $z=f(x,y)$ e inferiormente por $R$.
    \item El área de la base $R$.
    \item La distancia entre la gráfica y el plano $z=0$.
    \item La longitud del borde de $R$.
\end{multi}

\begin{multi}[]%
    {ID-Interp-01-v3}
    Si $f(x,y)\geq 0$, geométricamente la integral $\displaystyle\iint_R f\,dA$ mide:
    \item* El volumen bajo la superficie $z=f(x,y)$ sobre la región $R$.
    \item El área encerrada por la curva de nivel de $f$.
    \item El gradiente de $f$ en $R$.
    \item La superficie de la gráfica de $f$.
\end{multi}

\begin{multi}[]%
    {ID-Interp-02}
    La integral $\displaystyle\iint_R 1\,dA$ representa:
    \item* El área de la región $R$.
    \item El volumen bajo la gráfica de $f$.
    \item El perímetro de la región $R$.
    \item Siempre el valor $1$.
\end{multi}

\begin{multi}[]%
    {ID-Interp-02-v1}
    Al integrar la función constante $1$ sobre una región $R$, es decir $\displaystyle\iint_R 1\,dA$, se obtiene:
    \item* El área de $R$.
    \item El volumen de $R$.
    \item La longitud de la frontera de $R$.
    \item El número $1$.
\end{multi}

\begin{multi}[]%
    {ID-Interp-02-v2}
    ¿Qué cantidad geométrica se calcula con $\displaystyle\iint_R dA$?
    \item* El área de la región $R$.
    \item El perímetro de $R$.
    \item El volumen del sólido sobre $R$.
    \item El valor promedio de $R$.
\end{multi}

\begin{multi}[]%
    {ID-Interp-02-v3}
    Para hallar el área de una región plana $R$ mediante una integral doble se calcula:
    \item* $\displaystyle\iint_R 1\,dA$
    \item $\displaystyle\iint_R (x+y)\,dA$
    \item $\displaystyle\oint_{\partial R} 1\,ds$
    \item $\displaystyle\iint_R 0\,dA$
\end{multi}

\begin{multi}[]%
    {ID-Elem-01}
    En la expresión $\displaystyle\iint_R f(x,y)\,dA$, el símbolo $dA$ corresponde a:
    \item* El diferencial de área de la región de integración.
    \item El diferencial de longitud de arco.
    \item El diferencial de volumen.
    \item La función integrando.
\end{multi}

\begin{multi}[]%
    {ID-Elem-01-v1}
    ¿Qué representa $dA$ dentro de la integral doble $\displaystyle\iint_R f(x,y)\,dA$?
    \item* El elemento diferencial de área.
    \item El elemento diferencial de volumen.
    \item El diferencial de arco $ds$.
    \item El valor de $f$ en cada punto.
\end{multi}

\begin{multi}[]%
    {ID-Elem-01-v2}
    En coordenadas cartesianas, el diferencial de área $dA$ de una integral doble se escribe como:
    \item* $dx\,dy$
    \item $dx+dy$
    \item $dx\,dy\,dz$
    \item $\sqrt{dx^2+dy^2}$
\end{multi}

\begin{multi}[]%
    {ID-Elem-01-v3}
    El término $dA$ en $\displaystyle\iint_R f\,dA$ indica que se está acumulando $f$ respecto a:
    \item* Pequeños elementos de área de $R$.
    \item Pequeños elementos de longitud de $R$.
    \item Pequeños elementos de volumen.
    \item El tiempo.
\end{multi}

\begin{multi}[]%
    {ID-Fubini-01}
    El teorema de Fubini permite:
    \item* Calcular una integral doble mediante integrales iteradas (una variable a la vez).
    \item Intercambiar la integral por una derivada parcial.
    \item Transformar la integral doble en una integral de línea.
    \item Calcular el área únicamente en coordenadas polares.
\end{multi}

\begin{multi}[]%
    {ID-Fubini-01-v1}
    ¿Para qué sirve el teorema de Fubini?
    \item* Para reducir una integral doble a dos integrales simples sucesivas.
    \item Para convertir una integral doble en una integral triple.
    \item Para derivar bajo el signo integral.
    \item Para calcular el rotacional de un campo.
\end{multi}

\begin{multi}[]%
    {ID-Fubini-01-v2}
    El teorema de Fubini establece que, bajo condiciones de continuidad, una integral doble se puede evaluar:
    \item* Como integrales iteradas, integrando primero respecto a una variable y luego respecto a la otra.
    \item Únicamente sumando los valores de $f$ en la frontera.
    \item Solo después de un cambio a coordenadas polares.
    \item Derivando respecto a cada variable por separado.
\end{multi}

\begin{multi}[]%
    {ID-Fubini-01-v3}
    Gracias al teorema de Fubini, calcular $\displaystyle\iint_R f\,dA$ se reduce a:
    \item* Resolver una integral simple dentro de otra integral simple.
    \item Resolver una sola integral de una variable.
    \item Calcular un determinante jacobiano.
    \item Aplicar la regla de la cadena.
\end{multi}

\begin{multi}[]%
    {ID-Fubini-02}
    Si $R=[a,b]\times[c,d]$ es un rectángulo y $f$ es continua, el teorema de Fubini afirma que:
    \item* $\displaystyle\iint_R f\,dA = \int_a^b\!\!\int_c^d f(x,y)\,dy\,dx = \int_c^d\!\!\int_a^b f(x,y)\,dx\,dy$
    \item $\displaystyle\iint_R f\,dA = \int_a^b f(x,y)\,dx$
    \item $\displaystyle\iint_R f\,dA = \int_a^b\!\!\int_c^d f(x,y)\,dy\,dx$ solo en ese orden.
    \item $\displaystyle\iint_R f\,dA = \left(\int_a^b f\,dx\right)\left(\int_c^d f\,dy\right)$ siempre.
\end{multi}

\begin{multi}[]%
    {ID-Fubini-02-v1}
    Sobre un rectángulo $R=[a,b]\times[c,d]$ con $f$ continua, el orden de integración en las integrales iteradas:
    \item* Puede intercambiarse sin alterar el resultado.
    \item Debe ser siempre $dy\,dx$.
    \item Debe ser siempre $dx\,dy$.
    \item Cambia el signo del resultado al intercambiarse.
\end{multi}

\begin{multi}[]%
    {ID-Fubini-02-v2}
    Para $f$ continua sobre el rectángulo $R=[a,b]\times[c,d]$, ¿cuál igualdad es correcta?
    \item* $\displaystyle\int_a^b\!\!\int_c^d f\,dy\,dx = \int_c^d\!\!\int_a^b f\,dx\,dy$
    \item $\displaystyle\iint_R f\,dA = \int_a^b f\,dx + \int_c^d f\,dy$
    \item $\displaystyle\iint_R f\,dA = \left(\int_a^b f\,dx\right)\left(\int_c^d f\,dy\right)$
    \item $\displaystyle\iint_R f\,dA = \int_a^d f\,dx$
\end{multi}

\begin{multi}[]%
    {ID-Fubini-02-v3}
    En un rectángulo $R=[a,b]\times[c,d]$ con $f$ continua, el teorema de Fubini garantiza que las dos integrales iteradas posibles:
    \item* Dan el mismo valor.
    \item Difieren en un factor jacobiano.
    \item Solo coinciden si $f$ es separable.
    \item Tienen signos opuestos.
\end{multi}

\end{quiz}


%%%%%%%%%%%%%%%%%%%%%%%%%%%%%%%%%%%%%%%%
\begin{quiz}{Propiedades de la integral doble}
%%%%%%%%%%%%%%%%%%%%%%%%%%%%%%%%%%%%%%%%

\begin{multi}[]%
    {IDP-Lin-01}
    La propiedad de linealidad de la integral doble establece que:
    \item* $\displaystyle\iint_R (af+bg)\,dA = a\iint_R f\,dA + b\iint_R g\,dA$
    \item $\displaystyle\iint_R (f\cdot g)\,dA = \iint_R f\,dA \cdot \iint_R g\,dA$
    \item $\displaystyle\iint_R (f+g)\,dA = \iint_R f\,dA \cdot \iint_R g\,dA$
    \item $\displaystyle\iint_R (af)\,dA = \iint_R f\,dA$
\end{multi}

\begin{multi}[]%
    {IDP-Lin-01-v1}
    ¿Cuál de las siguientes igualdades expresa la linealidad de la integral doble?
    \item* $\displaystyle\iint_R (af+bg)\,dA = a\iint_R f\,dA + b\iint_R g\,dA$
    \item $\displaystyle\iint_R (af+bg)\,dA = ab\iint_R (f+g)\,dA$
    \item $\displaystyle\iint_R (af+bg)\,dA = \iint_R f\,dA + \iint_R g\,dA$
    \item $\displaystyle\iint_R (af+bg)\,dA = a+b$
\end{multi}

\begin{multi}[]%
    {IDP-Lin-01-v2}
    Una constante $c$ que multiplica al integrando de una integral doble:
    \item* Puede salir fuera de la integral: $\displaystyle\iint_R c\,f\,dA = c\iint_R f\,dA$.
    \item Debe elevarse al cuadrado al salir.
    \item No puede extraerse de la integral.
    \item Se convierte en el área de $R$.
\end{multi}

\begin{multi}[]%
    {IDP-Lin-01-v3}
    La integral doble de una suma de funciones cumple:
    \item* $\displaystyle\iint_R (f+g)\,dA = \iint_R f\,dA + \iint_R g\,dA$
    \item $\displaystyle\iint_R (f+g)\,dA = \iint_R f\,dA \cdot \iint_R g\,dA$
    \item $\displaystyle\iint_R (f+g)\,dA = \iint_R f\,dA - \iint_R g\,dA$
    \item $\displaystyle\iint_R (f+g)\,dA = \iint_R fg\,dA$
\end{multi}

\begin{multi}[]%
    {IDP-Adit-01}
    Si $R = R_1 \cup R_2$, donde $R_1$ y $R_2$ no se solapan (salvo en su frontera), entonces:
    \item* $\displaystyle\iint_R f\,dA = \iint_{R_1} f\,dA + \iint_{R_2} f\,dA$
    \item $\displaystyle\iint_R f\,dA = \iint_{R_1} f\,dA \cdot \iint_{R_2} f\,dA$
    \item $\displaystyle\iint_R f\,dA = \iint_{R_1} f\,dA - \iint_{R_2} f\,dA$
    \item $\displaystyle\iint_R f\,dA = \max\left(\iint_{R_1} f\,dA,\ \iint_{R_2} f\,dA\right)$
\end{multi}

\begin{multi}[]%
    {IDP-Adit-01-v1}
    La propiedad de aditividad respecto a la región (descomponer $R$ en dos partes que no se solapan) afirma que la integral sobre $R$ es:
    \item* La suma de las integrales sobre cada parte.
    \item El producto de las integrales sobre cada parte.
    \item El promedio de las integrales sobre cada parte.
    \item La mayor de las integrales sobre cada parte.
\end{multi}

\begin{multi}[]%
    {IDP-Adit-01-v2}
    Si una región $R$ se parte en subregiones $R_1$ y $R_2$ que solo comparten la frontera, $\displaystyle\iint_R f\,dA$ es igual a:
    \item* $\displaystyle\iint_{R_1} f\,dA + \iint_{R_2} f\,dA$
    \item $\displaystyle\iint_{R_1\cap R_2} f\,dA$
    \item $\displaystyle\iint_{R_1} f\,dA \cdot \iint_{R_2} f\,dA$
    \item $\dfrac{1}{2}\left(\displaystyle\iint_{R_1} f\,dA + \iint_{R_2} f\,dA\right)$
\end{multi}

\begin{multi}[]%
    {IDP-Adit-01-v3}
    Descomponer una región complicada en regiones más simples y sumar las integrales correspondientes es válido gracias a la propiedad de:
    \item* Aditividad respecto a la región de integración.
    \item Linealidad respecto al integrando.
    \item Monotonía.
    \item Cambio de variable.
\end{multi}

\begin{multi}[]%
    {IDP-Mon-01}
    Si $f(x,y)\leq g(x,y)$ para todo $(x,y)\in R$, entonces:
    \item* $\displaystyle\iint_R f\,dA \leq \iint_R g\,dA$
    \item $\displaystyle\iint_R f\,dA \geq \iint_R g\,dA$
    \item $\displaystyle\iint_R f\,dA = \iint_R g\,dA$
    \item No hay relación entre ambas integrales.
\end{multi}

\begin{multi}[]%
    {IDP-Mon-01-v1}
    La propiedad de monotonía de la integral doble afirma que si $f\leq g$ en toda la región $R$, entonces:
    \item* La integral de $f$ es menor o igual que la integral de $g$.
    \item La integral de $f$ es mayor que la de $g$.
    \item Ambas integrales son iguales.
    \item Las integrales tienen signos opuestos.
\end{multi}

\begin{multi}[]%
    {IDP-Mon-01-v2}
    Si $f(x,y)\geq 0$ para todo $(x,y)\in R$, ¿qué se puede afirmar de $\displaystyle\iint_R f\,dA$?
    \item* $\displaystyle\iint_R f\,dA \geq 0$
    \item \(\displaystyle\iint_R f\,dA < 0\)
    \item $\displaystyle\iint_R f\,dA = 0$ necesariamente.
    \item No puede determinarse su signo.
\end{multi}

\begin{multi}[]%
    {IDP-Mon-01-v3}
    Sabiendo que $m \leq f(x,y) \leq M$ en una región $R$ de área $A(R)$, la integral cumple:
    \item* $\displaystyle m\,A(R) \leq \iint_R f\,dA \leq M\,A(R)$
    \item $\displaystyle \iint_R f\,dA = M\,A(R)$
    \item $\displaystyle \iint_R f\,dA \geq M\,A(R)$
    \item $\displaystyle \iint_R f\,dA \leq m\,A(R)$
\end{multi}

\begin{multi}[]%
    {IDP-Cero-01}
    Si la región $R$ tiene área cero, entonces $\displaystyle\iint_R f\,dA$ es:
    \item* $0$
    \item $1$
    \item El valor de $f$ en el centro de $R$.
    \item Indefinida.
\end{multi}

\begin{multi}[]%
    {IDP-Cero-01-v1}
    El valor de $\displaystyle\iint_R f\,dA$ cuando $R$ es un conjunto de área nula (por ejemplo, un segmento o un punto) es:
    \item* $0$
    \item Igual al valor máximo de $f$.
    \item $1$
    \item Igual al área de la frontera.
\end{multi}

\begin{multi}[]%
    {IDP-Cero-01-v2}
    Si se integra la función nula $f\equiv 0$ sobre cualquier región $R$, el resultado es:
    \item* $0$
    \item El área de $R$.
    \item $1$
    \item Indefinido.
\end{multi}

\begin{multi}[]%
    {IDP-Cero-01-v3}
    ¿Cuánto vale la integral doble de cualquier función continua sobre una región de medida (área) cero?
    \item* Cero.
    \item El valor de $f$ en un punto.
    \item Uno.
    \item No está definida.
\end{multi}

\end{quiz}


%%%%%%%%%%%%%%%%%%%%%%%%%%%%%%%%%%%%%%%%
\begin{quiz}{Cambio de variable en integrales dobles}
%%%%%%%%%%%%%%%%%%%%%%%%%%%%%%%%%%%%%%%%

\begin{multi}[]%
    {CV-Jac-01}
    En el cambio de variable de una integral doble, el factor que corrige el diferencial de área es:
    \item* El valor absoluto del determinante del jacobiano de la transformación.
    \item El determinante de la matriz hessiana.
    \item El gradiente de la transformación.
    \item La traza del jacobiano.
\end{multi}

\begin{multi}[]%
    {CV-Jac-01-v1}
    Al aplicar un cambio de variable en una integral doble, el diferencial de área se multiplica por:
    \item* $\left|\det J\right|$, el valor absoluto del determinante del jacobiano.
    \item $\det H$, el determinante de la hessiana.
    \item La norma del gradiente.
    \item La derivada direccional de la transformación.
\end{multi}

\begin{multi}[]%
    {CV-Jac-01-v2}
    ¿Qué papel cumple el jacobiano en un cambio de variable de una integral doble?
    \item* Mide cómo la transformación escala las áreas.
    \item Mide la curvatura del integrando.
    \item Indica la dirección de máximo crecimiento.
    \item Determina la continuidad de $f$.
\end{multi}

\begin{multi}[]%
    {CV-Jac-01-v3}
    En la fórmula de cambio de variable $\displaystyle\iint_R f\,dA = \iint_{R'} f(\,\cdot\,)\,|\det J|\,du\,dv$, el término $|\det J|$ representa:
    \item* El factor de corrección del área entre los dos sistemas de coordenadas.
    \item El nuevo integrando.
    \item La región transformada.
    \item La densidad de $f$.
\end{multi}

\begin{multi}[]%
    {CV-Pol-01}
    Al pasar a coordenadas polares, el diferencial de área $dA = dx\,dy$ se transforma en:
    \item* $r\,dr\,d\theta$
    \item $dr\,d\theta$
    \item $r^2\,dr\,d\theta$
    \item $\dfrac{1}{r}\,dr\,d\theta$
\end{multi}

\begin{multi}[]%
    {CV-Pol-01-v1}
    En coordenadas polares, el elemento de área $dA$ se escribe como:
    \item* $r\,dr\,d\theta$
    \item $dr\,d\theta$
    \item $\theta\,dr\,d\theta$
    \item $r^2\,dr\,d\theta$
\end{multi}

\begin{multi}[]%
    {CV-Pol-01-v2}
    Al convertir $\displaystyle\iint_R f(x,y)\,dx\,dy$ a coordenadas polares, el factor adicional que aparece junto a $dr\,d\theta$ es:
    \item* $r$
    \item $r^2$
    \item $\dfrac{1}{r}$
    \item $\sin\theta$
\end{multi}

\begin{multi}[]%
    {CV-Pol-01-v3}
    El «$r$» extra que aparece en $dA = r\,dr\,d\theta$ al usar coordenadas polares proviene de:
    \item* El determinante del jacobiano de la transformación polar.
    \item La derivada de $f$ respecto a $r$.
    \item El radio de curvatura de la región.
    \item Una constante arbitraria.
\end{multi}

\begin{multi}[]%
    {CV-Pol-02}
    Las coordenadas polares resultan especialmente convenientes cuando la región de integración es:
    \item* Un círculo, un sector circular o un anillo.
    \item Un rectángulo con lados paralelos a los ejes.
    \item Un triángulo arbitrario.
    \item Cualquier región poligonal.
\end{multi}

\begin{multi}[]%
    {CV-Pol-02-v1}
    ¿En cuál de las siguientes regiones es más natural usar coordenadas polares?
    \item* Un disco centrado en el origen.
    \item Un cuadrado de lado $1$.
    \item Un trapecio.
    \item Una franja rectangular.
\end{multi}

\begin{multi}[]%
    {CV-Pol-02-v2}
    Las coordenadas polares simplifican una integral doble principalmente cuando la región o el integrando presentan:
    \item* Simetría circular (dependen de la distancia al origen).
    \item Lados paralelos a los ejes coordenados.
    \item Forma triangular.
    \item Discontinuidades aisladas.
\end{multi}

\begin{multi}[]%
    {CV-Pol-02-v3}
    Una integral sobre la región $x^2+y^2\leq 4$ se resuelve más fácilmente usando coordenadas:
    \item* Polares.
    \item Cartesianas.
    \item Esféricas.
    \item Cilíndricas.
\end{multi}

\begin{multi}[]%
    {CV-Jac-02}
    Para la transformación a coordenadas polares $x = r\cos\theta$, $y = r\sin\theta$, el determinante del jacobiano es:
    \item* $r$
    \item $1$
    \item $r^2$
    \item $\cos\theta + \sin\theta$
\end{multi}

\begin{multi}[]%
    {CV-Jac-02-v1}
    El jacobiano de la transformación $x = r\cos\theta$, $y = r\sin\theta$ vale:
    \item* $r$
    \item $\theta$
    \item $1$
    \item $r\cos\theta\sin\theta$
\end{multi}

\begin{multi}[]%
    {CV-Jac-02-v2}
    Al calcular $\displaystyle\frac{\partial(x,y)}{\partial(r,\theta)}$ para las coordenadas polares se obtiene:
    \item* $r$
    \item $-r$
    \item $r^2$
    \item $0$
\end{multi}

\begin{multi}[]%
    {CV-Jac-02-v3}
    El determinante de la matriz $\begin{pmatrix} \cos\theta & -r\sin\theta \\ \sin\theta & r\cos\theta \end{pmatrix}$, jacobiano de las coordenadas polares, es:
    \item* $r$
    \item $1$
    \item $r(\cos\theta-\sin\theta)$
    \item $r^2$
\end{multi}

\end{quiz}


%%%%%%%%%%%%%%%%%%%%%%%%%%%%%%%%%%%%%%%%
\begin{quiz}{Integral triple: concepto, elementos e interpretación}
%%%%%%%%%%%%%%%%%%%%%%%%%%%%%%%%%%%%%%%%

\begin{multi}[]%
    {IT-Concepto-01}
    Sea $f:\R^3\to\R$ continua sobre un sólido $W$. El resultado de $\displaystyle\iiint_W f\,dV$ es:
    \item* Un número real
    \item Un campo escalar
    \item Un campo vectorial
    \item Una superficie
\end{multi}

\begin{multi}[]%
    {IT-Concepto-01-v1}
    Al evaluar la integral triple $\displaystyle\iiint_W f\,dV$ de un campo escalar continuo sobre un sólido $W$, se obtiene:
    \item* Un número real
    \item Un campo vectorial
    \item Una función de tres variables
    \item Un sólido
\end{multi}

\begin{multi}[]%
    {IT-Concepto-01-v2}
    ¿Qué tipo de objeto resulta de calcular $\displaystyle\iiint_W f\,dV$ con $f:\R^3\to\R$ continua?
    \item* Un escalar (número real)
    \item Un vector de $\R^3$
    \item Una superficie en $\R^3$
    \item Una curva en el espacio
\end{multi}

\begin{multi}[]%
    {IT-Concepto-01-v3}
    El dominio de integración $W$ de una integral triple $\displaystyle\iiint_W f\,dV$ es:
    \item* Un sólido (región de $\R^3$).
    \item Una superficie de $\R^3$.
    \item Una curva de $\R^3$.
    \item Una región plana de $\R^2$.
\end{multi}

\begin{multi}[]%
    {IT-Interp-01}
    La integral $\displaystyle\iiint_W 1\,dV$ representa:
    \item* El volumen del sólido $W$.
    \item El área de la frontera de $W$.
    \item La masa total de $W$ con densidad arbitraria.
    \item El centro de masa de $W$.
\end{multi}

\begin{multi}[]%
    {IT-Interp-01-v1}
    Al integrar la función constante $1$ sobre un sólido $W$, es decir $\displaystyle\iiint_W 1\,dV$, se obtiene:
    \item* El volumen de $W$.
    \item El área de la superficie de $W$.
    \item La masa de $W$.
    \item El número $1$.
\end{multi}

\begin{multi}[]%
    {IT-Interp-01-v2}
    Para calcular el volumen de un sólido $W$ mediante una integral triple se evalúa:
    \item* $\displaystyle\iiint_W 1\,dV$
    \item $\displaystyle\iiint_W \rho\,dV$ con $\rho$ la densidad.
    \item $\displaystyle\iint_{\partial W} 1\,dS$
    \item $\displaystyle\iiint_W (x+y+z)\,dV$
\end{multi}

\begin{multi}[]%
    {IT-Interp-01-v3}
    ¿Qué cantidad geométrica calcula $\displaystyle\iiint_W dV$?
    \item* El volumen del sólido $W$.
    \item El área lateral de $W$.
    \item El perímetro de la base de $W$.
    \item La densidad media de $W$.
\end{multi}

\begin{multi}[]%
    {IT-Interp-02}
    Si $\rho(x,y,z)$ es la densidad de un sólido $W$, la integral $\displaystyle\iiint_W \rho\,dV$ representa:
    \item* La masa total del sólido.
    \item El volumen del sólido.
    \item La densidad promedio.
    \item El área de la superficie del sólido.
\end{multi}

\begin{multi}[]%
    {IT-Interp-02-v1}
    Para obtener la masa de un sólido $W$ con densidad puntual $\rho(x,y,z)$ se calcula:
    \item* $\displaystyle\iiint_W \rho\,dV$
    \item $\displaystyle\iiint_W 1\,dV$
    \item $\displaystyle\iiint_W \frac{1}{\rho}\,dV$
    \item $\displaystyle\iint_{\partial W} \rho\,dS$
\end{multi}

\begin{multi}[]%
    {IT-Interp-02-v2}
    La integral $\displaystyle\iiint_W \rho(x,y,z)\,dV$, donde $\rho$ es la densidad, da como resultado:
    \item* La masa total del sólido $W$.
    \item El volumen total de $W$.
    \item El peso de la superficie de $W$.
    \item El centro de masa de $W$.
\end{multi}

\begin{multi}[]%
    {IT-Interp-02-v3}
    Si la densidad de un sólido es constante $\rho_0$, su masa se obtiene como:
    \item* $\rho_0$ multiplicado por el volumen del sólido.
    \item $\rho_0$ multiplicado por el área de la superficie.
    \item El volumen dividido entre $\rho_0$.
    \item Simplemente $\rho_0$.
\end{multi}

\begin{multi}[]%
    {IT-Elem-01}
    En la expresión $\displaystyle\iiint_W f\,dV$, el símbolo $dV$ corresponde a:
    \item* El diferencial de volumen del sólido de integración.
    \item El diferencial de área.
    \item El diferencial de longitud de arco.
    \item La función integrando.
\end{multi}

\begin{multi}[]%
    {IT-Elem-01-v1}
    ¿Qué representa $dV$ en una integral triple $\displaystyle\iiint_W f\,dV$?
    \item* El elemento diferencial de volumen.
    \item El elemento diferencial de área.
    \item El diferencial de arco.
    \item El valor de $f$ en cada punto.
\end{multi}

\begin{multi}[]%
    {IT-Elem-01-v2}
    En coordenadas cartesianas, el diferencial de volumen $dV$ se escribe como:
    \item* $dx\,dy\,dz$
    \item $dx\,dy$
    \item $dx+dy+dz$
    \item $\sqrt{dx^2+dy^2+dz^2}$
\end{multi}

\begin{multi}[]%
    {IT-Elem-01-v3}
    El término $dV$ en $\displaystyle\iiint_W f\,dV$ indica que $f$ se acumula respecto a:
    \item* Pequeños elementos de volumen de $W$.
    \item Pequeños elementos de área de la frontera.
    \item Pequeños elementos de longitud.
    \item El número de variables.
\end{multi}

\begin{multi}[]%
    {IT-CilEsf-01}
    En coordenadas cilíndricas, el diferencial de volumen es:
    \item* $dV = r\,dr\,d\theta\,dz$
    \item $dV = dr\,d\theta\,dz$
    \item $dV = r^2\,dr\,d\theta\,dz$
    \item $dV = \rho^2\sin\varphi\,d\rho\,d\varphi\,d\theta$
\end{multi}

\begin{multi}[]%
    {IT-CilEsf-01-v1}
    El elemento de volumen en coordenadas cilíndricas se expresa como:
    \item* $r\,dr\,d\theta\,dz$
    \item $r^2\,dr\,d\theta\,dz$
    \item $dr\,d\theta\,dz$
    \item $\rho^2\,d\rho\,d\varphi\,d\theta$
\end{multi}

\begin{multi}[]%
    {IT-CilEsf-01-v2}
    Al pasar una integral triple a coordenadas cilíndricas, el factor jacobiano que multiplica a $dr\,d\theta\,dz$ es:
    \item* $r$
    \item $r^2$
    \item $\dfrac{1}{r}$
    \item $\sin\varphi$
\end{multi}

\begin{multi}[]%
    {IT-CilEsf-01-v3}
    En coordenadas cilíndricas $(r,\theta,z)$, el diferencial de volumen incluye el factor:
    \item* $r$ (el mismo de las coordenadas polares en el plano).
    \item $\rho^2\sin\varphi$.
    \item $z$.
    \item ninguno; es simplemente $dr\,d\theta\,dz$.
\end{multi}

\begin{multi}[]%
    {IT-CilEsf-02}
    En coordenadas esféricas, el diferencial de volumen es:
    \item* $dV = \rho^2\sin\varphi\,d\rho\,d\varphi\,d\theta$
    \item $dV = \rho\,d\rho\,d\varphi\,d\theta$
    \item $dV = r\,dr\,d\theta\,dz$
    \item $dV = \rho^2\,d\rho\,d\varphi\,d\theta$
\end{multi}

\begin{multi}[]%
    {IT-CilEsf-02-v1}
    El elemento de volumen en coordenadas esféricas $(\rho,\varphi,\theta)$ es:
    \item* $\rho^2\sin\varphi\,d\rho\,d\varphi\,d\theta$
    \item $\rho\sin\varphi\,d\rho\,d\varphi\,d\theta$
    \item $\rho^2\,d\rho\,d\varphi\,d\theta$
    \item $r\,dr\,d\theta\,dz$
\end{multi}

\begin{multi}[]%
    {IT-CilEsf-02-v2}
    El factor jacobiano que aparece al transformar una integral triple a coordenadas esféricas es:
    \item* $\rho^2\sin\varphi$
    \item $\rho\sin\varphi$
    \item $\rho^2$
    \item $r$
\end{multi}

\begin{multi}[]%
    {IT-CilEsf-02-v3}
    En coordenadas esféricas, ¿cuál es el factor que corrige el diferencial de volumen respecto de $d\rho\,d\varphi\,d\theta$?
    \item* $\rho^2\sin\varphi$
    \item $\sin\varphi$
    \item $\rho^2$
    \item $\cos\varphi$
\end{multi}

\begin{multi}[]%
    {IT-CilEsf-03}
    Para un sólido con simetría respecto a un eje (como un cilindro), el sistema de coordenadas más conveniente suele ser:
    \item* Cilíndricas
    \item Esféricas
    \item Cartesianas
    \item Polares en el plano
\end{multi}

\begin{multi}[]%
    {IT-CilEsf-03-v1}
    ¿Qué sistema de coordenadas conviene usar para integrar sobre un cilindro o un cono?
    \item* Coordenadas cilíndricas.
    \item Coordenadas esféricas.
    \item Coordenadas cartesianas.
    \item Coordenadas polares planas.
\end{multi}

\begin{multi}[]%
    {IT-CilEsf-03-v2}
    Un sólido invariante bajo rotaciones alrededor del eje $z$ se integra más cómodamente en coordenadas:
    \item* Cilíndricas.
    \item Esféricas.
    \item Cartesianas.
    \item Baricéntricas.
\end{multi}

\begin{multi}[]%
    {IT-CilEsf-03-v3}
    Para calcular el volumen de un sólido limitado por el cilindro $x^2+y^2=4$, el sistema más adecuado es:
    \item* Cilíndricas.
    \item Esféricas.
    \item Cartesianas exclusivamente.
    \item Polares en una sola variable.
\end{multi}

\begin{multi}[]%
    {IT-CilEsf-04}
    Para un sólido con simetría respecto a un punto (como una esfera o una bola), el sistema de coordenadas más conveniente suele ser:
    \item* Esféricas
    \item Cilíndricas
    \item Cartesianas
    \item Rectangulares
\end{multi}

\begin{multi}[]%
    {IT-CilEsf-04-v1}
    ¿Qué coordenadas conviene emplear para integrar sobre una bola centrada en el origen?
    \item* Coordenadas esféricas.
    \item Coordenadas cilíndricas.
    \item Coordenadas cartesianas.
    \item Coordenadas polares planas.
\end{multi}

\begin{multi}[]%
    {IT-CilEsf-04-v2}
    Un sólido con simetría radial respecto a un punto (depende solo de la distancia al origen) se integra más fácilmente en coordenadas:
    \item* Esféricas.
    \item Cilíndricas.
    \item Cartesianas.
    \item Rectangulares.
\end{multi}

\begin{multi}[]%
    {IT-CilEsf-04-v3}
    Para hallar el volumen de la región $x^2+y^2+z^2\leq 9$, el sistema de coordenadas más adecuado es:
    \item* Esféricas.
    \item Cilíndricas.
    \item Cartesianas únicamente.
    \item Polares en el plano.
\end{multi}

\end{quiz}


%%%%%%%%%%%%%%%%%%%%%%%%%%%%%%%%%%%%%%%%
\begin{quiz}{Integral de línea: concepto y elementos}
%%%%%%%%%%%%%%%%%%%%%%%%%%%%%%%%%%%%%%%%

\begin{multi}[]%
    {IL-Esc-01}
    La integral de línea de un campo escalar $f$ sobre una curva $C$, $\displaystyle\int_C f\,ds$, puede interpretarse físicamente como:
    \item* La masa de un alambre con forma $C$ y densidad lineal $f$.
    \item El trabajo realizado por una fuerza a lo largo de $C$.
    \item El flujo a través de $C$.
    \item El área encerrada por $C$.
\end{multi}

\begin{multi}[]%
    {IL-Esc-01-v1}
    Si $f$ representa la densidad lineal de un alambre con forma de la curva $C$, entonces $\displaystyle\int_C f\,ds$ corresponde a:
    \item* La masa total del alambre.
    \item El trabajo sobre el alambre.
    \item La longitud del alambre.
    \item El centro de masa del alambre.
\end{multi}

\begin{multi}[]%
    {IL-Esc-01-v2}
    ¿Qué cantidad física calcula la integral de línea de un campo escalar $\displaystyle\int_C f\,ds$?
    \item* La masa de un alambre de densidad $f$ y forma $C$.
    \item El trabajo de un campo vectorial.
    \item El flujo de un campo a través de $C$.
    \item La circulación del campo.
\end{multi}

\begin{multi}[]%
    {IL-Esc-01-v3}
    La integral $\displaystyle\int_C 1\,ds$ de un campo escalar constante igual a $1$ sobre la curva $C$ representa:
    \item* La longitud de la curva $C$.
    \item El área encerrada por $C$.
    \item El trabajo a lo largo de $C$.
    \item La masa de la región interior a $C$.
\end{multi}

\begin{multi}[]%
    {IL-Vec-01}
    La integral de línea de un campo vectorial $\displaystyle\int_C \vec{F}\cdot d\vec{r}$ se interpreta físicamente como:
    \item* El trabajo realizado por el campo $\vec{F}$ a lo largo de la curva $C$.
    \item La masa de un alambre con forma $C$.
    \item El volumen bajo la curva $C$.
    \item La longitud de la curva $C$.
\end{multi}

\begin{multi}[]%
    {IL-Vec-01-v1}
    ¿Qué magnitud física representa $\displaystyle\int_C \vec{F}\cdot d\vec{r}$ cuando $\vec{F}$ es un campo de fuerzas?
    \item* El trabajo efectuado por $\vec{F}$ al recorrer $C$.
    \item La masa del alambre $C$.
    \item El flujo de $\vec{F}$ a través de una superficie.
    \item La longitud de arco de $C$.
\end{multi}

\begin{multi}[]%
    {IL-Vec-01-v2}
    La integral de línea $\displaystyle\int_C \vec{F}\cdot d\vec{r}$ acumula a lo largo de la curva:
    \item* La componente tangencial del campo $\vec{F}$ (su trabajo).
    \item La componente normal del campo.
    \item La magnitud del campo únicamente.
    \item La densidad del medio.
\end{multi}

\begin{multi}[]%
    {IL-Vec-01-v3}
    Para una fuerza $\vec{F}$ que mueve una partícula a lo largo de una trayectoria $C$, el trabajo realizado se calcula con:
    \item* $\displaystyle\int_C \vec{F}\cdot d\vec{r}$
    \item $\displaystyle\int_C \|\vec{F}\|\,ds$
    \item $\displaystyle\int_C \vec{F}\,ds$
    \item $\displaystyle\iint_C \vec{F}\cdot d\vec{r}$
\end{multi}

\begin{multi}[]%
    {IL-Elem-01}
    En la integral de línea de un campo escalar, el elemento $ds$ representa:
    \item* El diferencial de longitud de arco de la curva.
    \item El diferencial de área.
    \item El vector tangente unitario.
    \item El diferencial de volumen.
\end{multi}

\begin{multi}[]%
    {IL-Elem-01-v1}
    ¿Qué significa $ds$ en la integral $\displaystyle\int_C f\,ds$?
    \item* Un diferencial de longitud de arco a lo largo de $C$.
    \item Un diferencial de área de la región interior.
    \item El diferencial del parámetro $t$.
    \item El vector tangente a $C$.
\end{multi}

\begin{multi}[]%
    {IL-Elem-01-v2}
    Si $C$ está parametrizada por $\vec{r}(t)$, el diferencial de longitud de arco $ds$ se expresa como:
    \item* $\|\vec{r}\,'(t)\|\,dt$
    \item $\vec{r}\,'(t)\,dt$
    \item $\vec{r}(t)\,dt$
    \item $dt$
\end{multi}

\begin{multi}[]%
    {IL-Elem-01-v3}
    El elemento $ds$ de una integral de línea escalar es siempre:
    \item* Una cantidad positiva (longitud).
    \item Un vector.
    \item Una cantidad que puede ser negativa.
    \item Un área.
\end{multi}

\begin{multi}[]%
    {IL-Param-01}
    Para calcular una integral de línea es necesario contar con:
    \item* Una parametrización de la curva $C$.
    \item El gradiente de la curva $C$.
    \item La matriz hessiana del integrando.
    \item El área de la región encerrada por $C$.
\end{multi}

\begin{multi}[]%
    {IL-Param-01-v1}
    El primer paso para evaluar una integral de línea sobre $C$ consiste en:
    \item* Parametrizar la curva $C$ mediante $\vec{r}(t)$.
    \item Calcular el área que encierra $C$.
    \item Hallar la hessiana del integrando.
    \item Determinar el rotacional del campo.
\end{multi}

\begin{multi}[]%
    {IL-Param-01-v2}
    Para transformar $\displaystyle\int_C \vec{F}\cdot d\vec{r}$ en una integral de una sola variable se requiere:
    \item* Una parametrización $\vec{r}(t)$, $t\in[a,b]$, de la curva.
    \item El gradiente del campo $\vec{F}$.
    \item La divergencia de $\vec{F}$.
    \item El área de $C$.
\end{multi}

\begin{multi}[]%
    {IL-Param-01-v3}
    ¿Qué elemento es indispensable para reducir una integral de línea a una integral ordinaria respecto al parámetro?
    \item* Una representación paramétrica de la curva.
    \item El valor del campo en el origen.
    \item La matriz jacobiana del campo.
    \item La longitud total de la curva.
\end{multi}

\begin{multi}[]%
    {IL-Orient-01}
    Si se invierte la orientación de la curva $C$, la integral de línea de un campo vectorial $\displaystyle\int_C \vec{F}\cdot d\vec{r}$:
    \item* Cambia de signo.
    \item No se altera.
    \item Se anula.
    \item Se duplica.
\end{multi}

\begin{multi}[]%
    {IL-Orient-01-v1}
    Al recorrer la curva $C$ en sentido contrario, el valor de $\displaystyle\int_C \vec{F}\cdot d\vec{r}$:
    \item* Cambia de signo.
    \item Permanece igual.
    \item Se hace cero.
    \item Se reduce a la mitad.
\end{multi}

\begin{multi}[]%
    {IL-Orient-01-v2}
    Si $-C$ denota la curva $C$ con orientación invertida, entonces:
    \item* $\displaystyle\int_{-C} \vec{F}\cdot d\vec{r} = -\int_{C} \vec{F}\cdot d\vec{r}$
    \item $\displaystyle\int_{-C} \vec{F}\cdot d\vec{r} = \int_{C} \vec{F}\cdot d\vec{r}$
    \item $\displaystyle\int_{-C} \vec{F}\cdot d\vec{r} = 0$
    \item $\displaystyle\int_{-C} \vec{F}\cdot d\vec{r} = 2\int_{C} \vec{F}\cdot d\vec{r}$
\end{multi}

\begin{multi}[]%
    {IL-Orient-01-v3}
    La integral de línea de un campo vectorial depende de la orientación de la curva porque:
    \item* El vector $d\vec{r}$ invierte su sentido al cambiar la orientación.
    \item El elemento $ds$ se vuelve negativo.
    \item El campo $\vec{F}$ cambia de magnitud.
    \item La parametrización deja de existir.
\end{multi}

\begin{multi}[]%
    {IL-Orient-02}
    La integral de línea de un campo escalar $\displaystyle\int_C f\,ds$ respecto a la orientación de la curva:
    \item* No depende de la orientación (es la misma en ambos sentidos).
    \item Cambia de signo al invertir la orientación.
    \item Se anula si la curva es cerrada.
    \item Depende solo del punto inicial.
\end{multi}

\begin{multi}[]%
    {IL-Orient-02-v1}
    Al recorrer la curva $C$ en sentido opuesto, la integral de línea escalar $\displaystyle\int_C f\,ds$:
    \item* Conserva su valor.
    \item Cambia de signo.
    \item Se anula.
    \item Se duplica.
\end{multi}

\begin{multi}[]%
    {IL-Orient-02-v2}
    A diferencia de la integral de línea vectorial, la integral de línea de un campo escalar:
    \item* Es independiente de la orientación de la curva.
    \item Cambia de signo al invertir la orientación.
    \item Solo está definida para curvas cerradas.
    \item Depende del jacobiano.
\end{multi}

\begin{multi}[]%
    {IL-Orient-02-v3}
    Si $-C$ es la curva $C$ con orientación invertida, para un campo escalar $f$ se cumple:
    \item* $\displaystyle\int_{-C} f\,ds = \int_{C} f\,ds$
    \item $\displaystyle\int_{-C} f\,ds = -\int_{C} f\,ds$
    \item $\displaystyle\int_{-C} f\,ds = 0$
    \item $\displaystyle\int_{-C} f\,ds = 2\int_{C} f\,ds$
\end{multi}

\end{quiz}


%%%%%%%%%%%%%%%%%%%%%%%%%%%%%%%%%%%%%%%%
\begin{quiz}{Campos conservativos}
%%%%%%%%%%%%%%%%%%%%%%%%%%%%%%%%%%%%%%%%

\begin{multi}[]%
    {CC-Def-01}
    Un campo vectorial $\vec{F}$ es conservativo cuando:
    \item* Proviene del gradiente de algún campo escalar, es decir, $\vec{F} = \nabla f$.
    \item Su divergencia es nula.
    \item Es perpendicular a sus curvas de nivel.
    \item Su jacobiano es invertible.
\end{multi}

\begin{multi}[]%
    {CC-Def-01-v1}
    ¿Qué caracteriza a un campo vectorial conservativo?
    \item* Existe un campo escalar $f$ tal que $\vec{F}=\nabla f$.
    \item Su divergencia es cero.
    \item Su rotacional es distinto de cero.
    \item Es constante en todo el dominio.
\end{multi}

\begin{multi}[]%
    {CC-Def-01-v2}
    Un campo $\vec{F}$ se llama conservativo (o campo de gradiente) si:
    \item* Admite una función potencial $f$ con $\vec{F}=\nabla f$.
    \item No admite función potencial.
    \item Es tangente a sus líneas de campo.
    \item Tiene divergencia constante.
\end{multi}

\begin{multi}[]%
    {CC-Def-01-v3}
    A la función escalar $f$ que satisface $\vec{F}=\nabla f$ se le denomina:
    \item* Función potencial de $\vec{F}$.
    \item Divergencia de $\vec{F}$.
    \item Rotacional de $\vec{F}$.
    \item Jacobiano de $\vec{F}$.
\end{multi}

\begin{multi}[]%
    {CC-Indep-01}
    En un campo vectorial conservativo, la integral de línea $\displaystyle\int_C \vec{F}\cdot d\vec{r}$:
    \item* No depende de la trayectoria, solo de los puntos inicial y final.
    \item Depende de la longitud de la curva $C$.
    \item Es siempre positiva.
    \item Cambia con cada parametrización de $C$.
\end{multi}

\begin{multi}[]%
    {CC-Indep-01-v1}
    Una consecuencia de que $\vec{F}$ sea conservativo es que su integral de línea:
    \item* Es independiente del camino seguido entre dos puntos.
    \item Depende fuertemente del camino elegido.
    \item Siempre vale cero.
    \item Solo se puede calcular en curvas cerradas.
\end{multi}

\begin{multi}[]%
    {CC-Indep-01-v2}
    Si $\vec{F}=\nabla f$, el teorema fundamental de las integrales de línea afirma que $\displaystyle\int_C \vec{F}\cdot d\vec{r}$ es igual a:
    \item* $f(B)-f(A)$, con $A$ y $B$ los extremos de $C$.
    \item La longitud de $C$.
    \item $f(A)+f(B)$.
    \item El área encerrada por $C$.
\end{multi}

\begin{multi}[]%
    {CC-Indep-01-v3}
    La propiedad de «independencia de la trayectoria» de las integrales de línea es equivalente a que el campo sea:
    \item* Conservativo.
    \item De divergencia nula.
    \item Constante.
    \item Perpendicular a sus líneas de flujo.
\end{multi}

\begin{multi}[]%
    {CC-Cerrada-01}
    Si $\vec{F}$ es un campo conservativo y $C$ es una curva cerrada, entonces $\displaystyle\oint_C \vec{F}\cdot d\vec{r}$ es:
    \item* $0$
    \item Siempre positiva.
    \item Igual a la longitud de $C$.
    \item Igual al área encerrada por $C$.
\end{multi}

\begin{multi}[]%
    {CC-Cerrada-01-v1}
    La circulación de un campo conservativo a lo largo de cualquier curva cerrada vale:
    \item* $0$
    \item $1$
    \item El área encerrada.
    \item El perímetro de la curva.
\end{multi}

\begin{multi}[]%
    {CC-Cerrada-01-v2}
    Si $\displaystyle\oint_C \vec{F}\cdot d\vec{r}=0$ para toda curva cerrada $C$, entonces el campo $\vec{F}$:
    \item* Es conservativo.
    \item Tiene rotacional no nulo.
    \item Depende del camino.
    \item No admite función potencial.
\end{multi}

\begin{multi}[]%
    {CC-Cerrada-01-v3}
    Para un campo conservativo, la integral de línea sobre un lazo cerrado se anula porque los puntos inicial y final:
    \item* Coinciden, de modo que $f(B)-f(A)=0$.
    \item Están muy alejados.
    \item Tienen la misma longitud de arco.
    \item Tienen distinto potencial.
\end{multi}

\begin{multi}[]%
    {CC-Grad-01}
    Decir que un campo vectorial $\vec{F}$ «proviene de un gradiente» significa que existe un campo escalar $f$ tal que:
    \item* $\vec{F} = \nabla f$
    \item $\vec{F} = \nabla \cdot f$
    \item $\vec{F} = \nabla \times f$
    \item $f = \nabla \vec{F}$
\end{multi}

\begin{multi}[]%
    {CC-Grad-01-v1}
    ¿Cuál expresión indica que $\vec{F}$ es el gradiente de un potencial escalar $f$?
    \item* $\vec{F} = \nabla f$
    \item $\vec{F} = \operatorname{div} f$
    \item $\vec{F} = \operatorname{rot} f$
    \item $\vec{F} = \Delta f$
\end{multi}

\begin{multi}[]%
    {CC-Grad-01-v2}
    Que $\vec{F}$ sea un «campo de gradiente» equivale a escribir:
    \item* $\vec{F} = \nabla f$ para alguna función escalar $f$.
    \item $\nabla\cdot\vec{F} = 0$.
    \item $\nabla\times\vec{F} = \vec{F}$.
    \item $f = \nabla\cdot\vec{F}$.
\end{multi}

\begin{multi}[]%
    {CC-Grad-01-v3}
    La operación $\nabla f$ que aparece en $\vec{F}=\nabla f$ corresponde a:
    \item* El gradiente del campo escalar $f$.
    \item La divergencia de $f$.
    \item El rotacional de $f$.
    \item El laplaciano de $f$.
\end{multi}

\begin{multi}[]%
    {CC-Grad-02}
    Si $\func{F}{\R^3}{\R^3}$ proviene de un gradiente (es conservativo), entonces necesariamente su rotacional cumple:
    \item* $\nabla \times \vec{F} = \vec{0}$
    \item $\nabla \times \vec{F} \neq \vec{0}$
    \item $\nabla \cdot \vec{F} = 0$
    \item $\nabla \cdot \vec{F} = 1$
\end{multi}

\begin{multi}[]%
    {CC-Grad-02-v1}
    Una condición necesaria para que un campo $\func{F}{\R^3}{\R^3}$ de clase $C^1$ sea conservativo es:
    \item* Que su rotacional sea nulo, $\nabla\times\vec{F}=\vec{0}$.
    \item Que su divergencia sea nula.
    \item Que su divergencia valga $1$.
    \item Que su rotacional sea no nulo.
\end{multi}

\begin{multi}[]%
    {CC-Grad-02-v2}
    Si un campo $\vec{F}$ en $\R^3$ admite función potencial, entonces $\nabla\times\vec{F}$ vale:
    \item* $\vec{0}$
    \item $\vec{F}$
    \item $1$
    \item $\nabla\cdot\vec{F}$
\end{multi}

\begin{multi}[]%
    {CC-Grad-02-v3}
    ¿Qué propiedad del rotacional sirve como prueba (necesaria) para descartar que un campo en $\R^3$ sea conservativo?
    \item* Si $\nabla\times\vec{F}\neq\vec{0}$, el campo no es conservativo.
    \item Si $\nabla\cdot\vec{F}\neq 0$, el campo no es conservativo.
    \item Si $\nabla\times\vec{F}=\vec{0}$, el campo no es conservativo.
    \item El rotacional no aporta información.
\end{multi}

\begin{multi}[]%
    {CC-Grad-03}
    ¿Cuál de las siguientes condiciones permite reconocer que un campo $\vec{F}=(P,Q)$ de clase $C^1$ podría provenir de un gradiente en un dominio simplemente conexo?
    \item* $\dfrac{\partial P}{\partial y} = \dfrac{\partial Q}{\partial x}$
    \item $\dfrac{\partial P}{\partial x} = \dfrac{\partial Q}{\partial y}$
    \item $P = Q$
    \item $\dfrac{\partial P}{\partial x} + \dfrac{\partial Q}{\partial y} = 0$
\end{multi}

\begin{multi}[]%
    {CC-Grad-03-v1}
    Para un campo plano $\vec{F}=(P,Q)$ de clase $C^1$ en un dominio simplemente conexo, la condición que garantiza que sea conservativo es:
    \item* $\dfrac{\partial P}{\partial y} = \dfrac{\partial Q}{\partial x}$
    \item $\dfrac{\partial P}{\partial x} = \dfrac{\partial Q}{\partial y}$
    \item $\dfrac{\partial P}{\partial y} = -\dfrac{\partial Q}{\partial x}$
    \item $P\cdot Q = 0$
\end{multi}

\begin{multi}[]%
    {CC-Grad-03-v2}
    La igualdad $\dfrac{\partial P}{\partial y} = \dfrac{\partial Q}{\partial x}$ para un campo $\vec{F}=(P,Q)$ en un dominio simplemente conexo indica que:
    \item* El campo es conservativo (admite función potencial).
    \item El campo tiene divergencia nula.
    \item El campo es solenoidal.
    \item El campo no tiene función potencial.
\end{multi}

\begin{multi}[]%
    {CC-Grad-03-v3}
    Para el campo plano $\vec{F}=(P,Q)$, la cantidad $\dfrac{\partial Q}{\partial x}-\dfrac{\partial P}{\partial y}$ corresponde a la componente del rotacional; el campo es conservativo (en un dominio simplemente conexo) cuando dicha cantidad es:
    \item* Igual a $0$.
    \item Positiva.
    \item Igual a $P+Q$.
    \item Distinta de $0$.
\end{multi}

\end{quiz}


%%%%%%%%%%%%%%%%%%%%%%%%%%%%%%%%%%%%%%%%
\begin{quiz}{Cálculo de integrales: evaluación numérica}
%%%%%%%%%%%%%%%%%%%%%%%%%%%%%%%%%%%%%%%%

\begin{numerical}[tolerance=0.01]%
    {IN-Calc-01}
    Sea $R = [0,2]\times[0,3]$. El valor de $\displaystyle\iint_R 1\,dA$ (el área de $R$) es:
    \item 6
\end{numerical}

\begin{numerical}[tolerance=0.01]%
    {IN-Calc-01-v1}
    Sea $R = [0,4]\times[0,2]$. El valor de $\displaystyle\iint_R 1\,dA$ (el área de $R$) es:
    \item 8
\end{numerical}

\begin{numerical}[tolerance=0.01]%
    {IN-Calc-01-v2}
    El área de la región rectangular $R=[1,4]\times[0,2]$, calculada como $\displaystyle\iint_R dA$, es:
    \item 6
\end{numerical}

\begin{numerical}[tolerance=0.01]%
    {IN-Calc-01-v3}
    Sea $R = [0,5]\times[0,2]$. ¿Cuánto vale $\displaystyle\iint_R 1\,dA$?
    \item 10
\end{numerical}

\begin{numerical}[tolerance=0.01]%
    {IN-Calc-02}
    Sea $R = [0,1]\times[0,1]$. El valor de $\displaystyle\iint_R (x+y)\,dA$ es:
    \item 1
\end{numerical}

\begin{numerical}[tolerance=0.01]%
    {IN-Calc-02-v1}
    Sea $R = [0,1]\times[0,1]$. El valor de $\displaystyle\iint_R x\,dA$ es:
    \item 0.5
\end{numerical}

\begin{numerical}[tolerance=0.01]%
    {IN-Calc-02-v2}
    Sea $R = [0,2]\times[0,1]$. El valor de $\displaystyle\iint_R (x+y)\,dA$ es:
    \item 3
\end{numerical}

\begin{numerical}[tolerance=0.01]%
    {IN-Calc-02-v3}
    Sea $R = [0,1]\times[0,2]$. El valor de $\displaystyle\iint_R (x+y)\,dA$ es:
    \item 3
\end{numerical}

\begin{numerical}[tolerance=0.01]%
    {IN-Calc-03}
    Sea $R = [0,2]\times[0,1]$. El valor de $\displaystyle\iint_R x\,dA$ es:
    \item 2
\end{numerical}

\begin{numerical}[tolerance=0.01]%
    {IN-Calc-03-v1}
    Sea $R = [0,3]\times[0,1]$. El valor de $\displaystyle\iint_R x\,dA$ es:
    \item 4.5
\end{numerical}

\begin{numerical}[tolerance=0.01]%
    {IN-Calc-03-v2}
    Sea $R = [0,2]\times[0,2]$. El valor de $\displaystyle\iint_R x\,dA$ es:
    \item 4
\end{numerical}

\begin{numerical}[tolerance=0.01]%
    {IN-Calc-03-v3}
    Sea $R = [0,2]\times[0,1]$. El valor de $\displaystyle\iint_R y\,dA$ es:
    \item 1
\end{numerical}

\begin{numerical}[tolerance=0.01]%
    {IN-Calc-04}
    El valor del determinante del jacobiano de la transformación a coordenadas polares $x=r\cos\theta$, $y=r\sin\theta$ evaluado en $r=3$ es:
    \item 3
\end{numerical}

\begin{numerical}[tolerance=0.01]%
    {IN-Calc-04-v1}
    El determinante del jacobiano de la transformación a coordenadas polares $x=r\cos\theta$, $y=r\sin\theta$ evaluado en $r=5$ es:
    \item 5
\end{numerical}

\begin{numerical}[tolerance=0.01]%
    {IN-Calc-04-v2}
    En coordenadas polares, el diferencial de área es $dA=r\,dr\,d\theta$. El valor del factor $r$ cuando $r=7$ es:
    \item 7
\end{numerical}

\begin{numerical}[tolerance=0.01]%
    {IN-Calc-04-v3}
    El determinante del jacobiano de las coordenadas polares evaluado en $r=1/2$ es:
    \item 0.5
\end{numerical}

\begin{numerical}[tolerance=0.01]%
    {IN-Calc-05}
    Sea $W = [0,1]\times[0,2]\times[0,3]$. El valor de $\displaystyle\iiint_W 1\,dV$ (el volumen de $W$) es:
    \item 6
\end{numerical}

\begin{numerical}[tolerance=0.01]%
    {IN-Calc-05-v1}
    Sea $W = [0,2]\times[0,2]\times[0,2]$. El valor de $\displaystyle\iiint_W 1\,dV$ (el volumen de $W$) es:
    \item 8
\end{numerical}

\begin{numerical}[tolerance=0.01]%
    {IN-Calc-05-v2}
    El volumen de la caja $W=[0,1]\times[0,1]\times[0,4]$, calculado como $\displaystyle\iiint_W dV$, es:
    \item 4
\end{numerical}

\begin{numerical}[tolerance=0.01]%
    {IN-Calc-05-v3}
    Sea $W = [0,3]\times[0,2]\times[0,2]$. ¿Cuánto vale $\displaystyle\iiint_W 1\,dV$?
    \item 12
\end{numerical}

\begin{numerical}[tolerance=0.01]%
    {IN-Calc-06}
    Sea $f(x,y,z) = x^2 + y^2$. En coordenadas cilíndricas, el integrando $f$ se escribe como $r^2$. El valor de $r^2$ cuando $r=2$ es:
    \item 4
\end{numerical}

\begin{numerical}[tolerance=0.01]%
    {IN-Calc-06-v1}
    Sea $f(x,y,z) = x^2 + y^2$. En coordenadas cilíndricas $f = r^2$. El valor de $r^2$ cuando $r=3$ es:
    \item 9
\end{numerical}

\begin{numerical}[tolerance=0.01]%
    {IN-Calc-06-v2}
    En coordenadas esféricas, $x^2+y^2+z^2 = \rho^2$. El valor de $\rho^2$ cuando $\rho=4$ es:
    \item 16
\end{numerical}

\begin{numerical}[tolerance=0.01]%
    {IN-Calc-06-v3}
    En coordenadas cilíndricas, el factor del diferencial de volumen es $r$. Si $r=2$, el valor de $dV/(dr\,d\theta\,dz)=r$ es:
    \item 2
\end{numerical}

\begin{numerical}[tolerance=0.01]%
    {IN-Calc-07}
    Sea $\vec{F} = \nabla f$ con $f(x,y) = x^2 + y^2$. El valor de la integral de línea $\displaystyle\int_C \vec{F}\cdot d\vec{r}$ a lo largo de cualquier curva $C$ que va de $(0,0)$ a $(1,2)$ es:
    \item 5
\end{numerical}

\begin{numerical}[tolerance=0.01]%
    {IN-Calc-07-v1}
    Sea $\vec{F} = \nabla f$ con $f(x,y) = x^2 + y^2$. El valor de $\displaystyle\int_C \vec{F}\cdot d\vec{r}$ a lo largo de cualquier curva que va de $(0,0)$ a $(2,1)$ es:
    \item 5
\end{numerical}

\begin{numerical}[tolerance=0.01]%
    {IN-Calc-07-v2}
    Sea $\vec{F} = \nabla f$ con $f(x,y) = x^2 + y^2$. El valor de $\displaystyle\int_C \vec{F}\cdot d\vec{r}$ desde $(0,0)$ hasta $(3,0)$ es:
    \item 9
\end{numerical}

\begin{numerical}[tolerance=0.01]%
    {IN-Calc-07-v3}
    Sea $\vec{F} = \nabla f$ con $f(x,y) = x^2 + y^2$. El valor de $\displaystyle\int_C \vec{F}\cdot d\vec{r}$ a lo largo de cualquier curva que va de $(1,0)$ a $(2,0)$ es:
    \item 3
\end{numerical}

\begin{numerical}[tolerance=0.01]%
    {IN-Calc-08}
    Sea $\vec{F} = \nabla f$ con $f(x,y) = xy$. El valor de $\displaystyle\int_C \vec{F}\cdot d\vec{r}$ a lo largo de una curva $C$ que va de $(1,1)$ a $(3,2)$ es:
    \item 5
\end{numerical}

\begin{numerical}[tolerance=0.01]%
    {IN-Calc-08-v1}
    Sea $\vec{F} = \nabla f$ con $f(x,y) = xy$. El valor de $\displaystyle\int_C \vec{F}\cdot d\vec{r}$ desde $(0,0)$ hasta $(2,3)$ es:
    \item 6
\end{numerical}

\begin{numerical}[tolerance=0.01]%
    {IN-Calc-08-v2}
    Sea $\vec{F} = \nabla f$ con $f(x,y) = xy$. El valor de $\displaystyle\int_C \vec{F}\cdot d\vec{r}$ desde $(1,1)$ hasta $(4,2)$ es:
    \item 7
\end{numerical}

\begin{numerical}[tolerance=0.01]%
    {IN-Calc-08-v3}
    Sea $\vec{F} = \nabla f$ con $f(x,y) = xy$. El valor de $\displaystyle\oint_C \vec{F}\cdot d\vec{r}$ a lo largo de cualquier curva cerrada $C$ es:
    \item 0
\end{numerical}

\end{quiz}

\end{document}
